\documentclass[11pt,a4paper]{article}
\usepackage[spanish,es-noquoting]{babel}
\usepackage[top=2.2cm, bottom=2.2cm, left=2.5cm, right=2.5cm]{geometry}
\usepackage{fontspec}
\usepackage{amsmath}
\usepackage[version=4]{mhchem}
\usepackage{enumitem}
\usepackage{microtype}

\setmainfont{Calibri}
\linespread{1.15}

\begin{document}
\pagestyle{empty}

\begin{center}
    {\LARGE \textbf{Ejercicios Básicos de Rendimiento de Reacciones}}\\[4pt]
    {\large \textbf{Física y Química 4.º ESO / 1.º Bachillerato}}
\end{center}

\vspace{0.4cm}
\noindent\textbf{1.} En el laboratorio se calientan $122{,}5\text{ g}$ de clorato de potasio (\ce{KClO3}) puro, produciéndose su descomposición térmica completa en cloruro de potasio (\ce{KCl}) y oxígeno gaseoso (\ce{O2}) según la siguiente ecuación ajustada:
\[
\ce{2 KClO3 (s) -> 2 KCl (s) + 3 O2 (g)}
\]

\begin{enumerate}[label=\textbf{\alph*)}, leftmargin=2em, itemsep=6pt]
    \item Calcula los moles y la masa teórica máxima de gas oxígeno (\ce{O2}) que deberían obtenerse si el rendimiento del proceso fuese del $100\%$.
    \item Si al finalizar la experiencia se recogen experimentalmente en el laboratorio $38{,}4\text{ g}$ de \ce{O2}, calcula el porcentaje de rendimiento ($\%$) de la reacción.
\end{enumerate}

\vspace{0.25cm}
\noindent\textbf{Datos:} Masas atómicas ($\text{u}$): $\mathrm{O} = 16$; $\mathrm{Cl} = 35{,}5$; $\mathrm{K} = 39$.

\vspace{0.7cm}
\noindent\textbf{2.} Se hacen reaccionar $65{,}4\text{ g}$ de zinc metálico (\ce{Zn}) puro con una cantidad suficiente de ácido clorhídrico (\ce{HCl}) según la reacción ajustada:
\[
\ce{Zn (s) + 2 HCl (aq) -> ZnCl2 (aq) + H2 (g)}
\]
Sabiendo que las pérdidas del proceso hacen que el rendimiento real de la reacción sea del $75\%$, determina:
\begin{enumerate}[label=\textbf{\alph*)}, leftmargin=2em, itemsep=6pt]
    \item La masa teórica de cloruro de zinc (\ce{ZnCl2}) esperada al $100\%$ de rendimiento.
    \item La masa real de cloruro de zinc (\ce{ZnCl2}) que se obtiene en la práctica.
    \item El volumen real de gas hidrógeno (\ce{H2}) desprendido, medido en Condiciones Normales ($\text{C.N.}$: $1\text{ atm}$, $0\,^\circ\text{C}$, donde $1\text{ mol} = 22{,}4\text{ L}$).
\end{enumerate}

\vspace{0.25cm}
\noindent\textbf{Datos:} Masas atómicas ($\text{u}$): $\mathrm{H} = 1$; $\mathrm{Cl} = 35{,}5$; $\mathrm{Zn} = 65{,}4$.

\end{document}
